% GENERATED FILE. Edit canonical lessons, then run npm run generate:books.
% canonical-source-hash: fnv1a64:6cf93ba937040a71
% canonical-lessons: SA-C137-sivyati, SA-C137-jayati, SA-C137-guhati, SA-C137-raksati, SA-C137-jvalati

\chapter{Sews, Wins, Hides, Protects, Burns}
\label{ch:sa-sews-wins-hides-protects-burns}
\hlchaptermodality{137}

\begin{chapteropening}
\textbf{By the end of this chapter:} \emph{I can say sews, wins, hides, protects and burns.}

Complete the last lesson of chapter 137: I can say sews, wins, hides, protects and burns.
\end{chapteropening}

\section[\texorpdfstring{\sk{सीव्यति} (sīvyati)}{सीव्यति (sīvyati)}]{\sk{सीव्यति} (sīvyati) — sews}

\label{lesson:SA-C137-sivyati}



\begin{quote}
Before the new one: say the Sanskrit for wipes, then the Sanskrit for ties.
\end{quote}

\subsection*{You'll want to know: \sk{सीव्यति}}
\textbf{\sk{सीव्यति}} — \emph{sīvyati} — \textquotedblleft{}sews\textquotedblright{}.

\begin{grammarlens}[title={What you've built}]
One more word for this chapter.
\end{grammarlens}

\subsection*{Your turn}
\begin{itemize}
\raggedright
  \item \emph{Say it:} \emph{sīvyati}
  \item \emph{Say it:} \emph{sīvyati}, once more
  \item \emph{Say it:} say \emph{sīvyati}
  \item \emph{Recall:} say the Sanskrit for wipes, then the Sanskrit for ties, then say \emph{sīvyati} again
\end{itemize}

\begin{tcolorbox}[breakable,colback=teal!4,colframe=teal!35!black,title={Before you move on}]
What does \sk{सीव्यति} mean? (\textquotedblleft{}Sews\textquotedblright{}.) Say it once more.
\end{tcolorbox}
\section[\texorpdfstring{\sk{जयति} (jayati)}{जयति (jayati)}]{\sk{जयति} (jayati) — wins}

\label{lesson:SA-C137-jayati}



\begin{quote}
Before the new one: say the Sanskrit for ties, then the Sanskrit for sews.
\end{quote}

\subsection*{You'll want to know: \sk{जयति}}
\textbf{\sk{जयति}} — \emph{jayati} — \textquotedblleft{}wins\textquotedblright{}.

\begin{grammarlens}[title={What you've built}]
One more word for this chapter.
\end{grammarlens}

\subsection*{Your turn}
\begin{itemize}
\raggedright
  \item \emph{Say it:} \emph{jayati}
  \item \emph{Say it:} \emph{jayati}, once more
  \item \emph{Say it:} say \emph{jayati}
  \item \emph{Recall:} say the Sanskrit for ties, then the Sanskrit for sews, then say \emph{jayati} again
\end{itemize}

\begin{tcolorbox}[breakable,colback=teal!4,colframe=teal!35!black,title={Before you move on}]
What does \sk{जयति} mean? (\textquotedblleft{}Wins\textquotedblright{}.) Say it once more.
\end{tcolorbox}
\section[\texorpdfstring{\sk{गूहति} (gūhati)}{गूहति (gūhati)}]{\sk{गूहति} (gūhati) — hides}

\label{lesson:SA-C137-guhati}



\begin{quote}
Before the new one: say the Sanskrit for sews, then the Sanskrit for wins.
\end{quote}

\subsection*{You'll want to know: \sk{गूहति}}
\textbf{\sk{गूहति}} — \emph{gūhati} — \textquotedblleft{}hides\textquotedblright{}.

\begin{grammarlens}[title={What you've built}]
One more word for this chapter.
\end{grammarlens}

\subsection*{Your turn}
\begin{itemize}
\raggedright
  \item \emph{Say it:} \emph{gūhati}
  \item \emph{Say it:} \emph{gūhati}, once more
  \item \emph{Say it:} say \emph{gūhati}
  \item \emph{Recall:} say the Sanskrit for sews, then the Sanskrit for wins, then say \emph{gūhati} again
\end{itemize}

\begin{tcolorbox}[breakable,colback=teal!4,colframe=teal!35!black,title={Before you move on}]
What does \sk{गूहति} mean? (\textquotedblleft{}Hides\textquotedblright{}.) Say it once more.
\end{tcolorbox}
\section[\texorpdfstring{\sk{रक्षति} (rakṣati)}{रक्षति (rakṣati)}]{\sk{रक्षति} (rakṣati) — protects}

\label{lesson:SA-C137-raksati}



\begin{quote}
Before the new one: say the Sanskrit for wins, then the Sanskrit for hides.
\end{quote}

\subsection*{You'll want to know: \sk{रक्षति}}
\textbf{\sk{रक्षति}} — \emph{rakṣati} — \textquotedblleft{}protects\textquotedblright{}.

\begin{grammarlens}[title={What you've built}]
One more word for this chapter.
\end{grammarlens}

\subsection*{Your turn}
\begin{itemize}
\raggedright
  \item \emph{Say it:} \emph{rakṣati}
  \item \emph{Say it:} \emph{rakṣati}, once more
  \item \emph{Say it:} say \emph{rakṣati}
  \item \emph{Recall:} say the Sanskrit for wins, then the Sanskrit for hides, then say \emph{rakṣati} again
\end{itemize}

\begin{tcolorbox}[breakable,colback=teal!4,colframe=teal!35!black,title={Before you move on}]
What does \sk{रक्षति} mean? (\textquotedblleft{}Protects\textquotedblright{}.) Say it once more.
\end{tcolorbox}
\section[\texorpdfstring{\sk{ज्वलति} (jvalati)}{ज्वलति (jvalati)}]{\sk{ज्वलति} (jvalati) — burns}

\label{lesson:SA-C137-jvalati}



\begin{quote}
Before the new one: say the Sanskrit for hides, then the Sanskrit for protects.
\end{quote}

\subsection*{You'll want to know: \sk{ज्वलति}}
\textbf{\sk{ज्वलति}} — \emph{jvalati} — \textquotedblleft{}burns\textquotedblright{}.

\begin{grammarlens}[title={What you've built}]
One more word for this chapter.
\end{grammarlens}

\subsection*{Your turn}
\begin{itemize}
\raggedright
  \item \emph{Say it:} \emph{jvalati}
  \item \emph{Say it:} \emph{jvalati}, once more
  \item \emph{Say it:} say \emph{jvalati}
  \item \emph{Recall:} say the Sanskrit for hides, then the Sanskrit for protects, then say \emph{jvalati} again
\end{itemize}

\begin{tcolorbox}[breakable,colback=teal!4,colframe=teal!35!black,title={Before you move on}]
What does \sk{ज्वलति} mean? (\textquotedblleft{}Burns\textquotedblright{}.) Say it once more.
\end{tcolorbox}
